\chapter{Prêt de titres et vente à découvert en pratique}\label{ch:m1:stock-loan-and-short-selling}

Un fonds est vendeur à découvert d'une action qu'il estime valoir la moitié de son cours, et paie 1~\% par
an pour l'emprunter. D'autres parviennent à la même conclusion~; les actions que les prêteurs proposent
s'épuisent~; en une semaine, les frais atteignent 80~\%. Le fonds paie désormais deux dixièmes de
pour cent de la position chaque jour pour conserver une position qui peut mettre un an à
porter ses fruits. Puis un prêteur vend ses actions et rappelle le prêt~; le courtier ne
trouve pas de titres de remplacement, et le fonds doit subir un \omterm{def:m1:stock-loan-and-short-selling:recall}{buy-in} sur un marché en hausse, avec
tous les autres vendeurs à découvert. Une position longue peut avoir tort pendant des années, sans autre coût que la patience.
Une position vendeuse est louée~: le loyer est fixé chaque jour par une enchère que le locataire
ne voit pas, et le bailleur peut mettre fin au bail du jour au lendemain. Le \Cref{ch:m1:financing}
a défini le prêt~; ce chapitre traite de la manière de vivre avec lui.

\section{La chaîne de prêt}

\begin{definition}[Offre prêtable et taux d'utilisation]\label{def:m1:stock-loan-and-short-selling:utilisation}
L'\emph{offre prêtable}\index{offre prêtable} d'un titre est la
quantité que ses détenteurs ont rendue disponible au prêt par l'intermédiaire de leurs agents.
Le \emph{taux d'utilisation}\index{taux d'utilisation} est la fraction de cette offre
actuellement en prêt.
\end{definition}

Les détenteurs sont des fonds de pension, des fonds indiciels, des assureurs et des fonds souverains~:
des détenteurs qui ne vendront pas la semaine prochaine et accueillent volontiers quelques points de base
de revenu. Ils prêtent par l'intermédiaire d'\emph{agents prêteurs}, en général leurs dépositaires, aux
\omterm{def:m1:financing:pb}{prime brokers} de la \cref{def:m1:financing:pb}, qui prêtent à leur tour à leurs
clients. Les \omterm{def:m1:financing:pb}{prime brokers} prêtent aussi à partir des comptes de marge de leurs clients
(\omterm{def:m1:financing:pb}{réhypothécation}) et s'empruntent entre eux. Les détenteurs longs qui n'ont pas
accepté de prêter (la plupart des comptes espèces de particuliers, de nombreux gérants actifs) ne figurent
pas du tout dans l'offre~: pour une grande valeur type, l'\omterm{def:m1:stock-loan-and-short-selling:utilisation}{offre prêtable} ne représente qu'une
fraction du flottant, et ce qui est réellement prêté n'en représente qu'une petite
fraction.

\begin{omfigure}
\begin{tikzpicture}[font=\small,
  box/.style={draw=omDef,rounded corners=2pt,minimum width=2.45cm,minimum height=1.05cm,align=center,fill=omDef!4},
  flow/.style={-{Stealth[length=2mm]},thick}]
\node[box] (own) at (0,0) {Bénéficiaire\\effectif};
\node[box] (lend) at (3.45,0) {Agent\\prêteur};
\node[box] (pb) at (6.9,0) {Prime\\broker};
\node[box] (hf) at (10.35,0) {Vendeur\\à découvert};
\node[box,draw=omExo,fill=omExo!6] (buyer) at (10.35,-2.4) {Acheteur\\sur le marché};
\foreach \a/\b in {own/lend,lend/pb,pb/hf}{
  \draw[flow,omDef] ([yshift=7pt]\a.east) -- ([yshift=7pt]\b.west);
  \draw[flow,omThm] ([yshift=-7pt]\b.west) -- ([yshift=-7pt]\a.east);}
\draw[flow,omDef] (hf.south) -- node[right,font=\footnotesize]{vend les actions} (buyer.north);
\node[font=\footnotesize,text=omDef] at (5.2,1.0) {actions (transfert de propriété~; rappelables à tout moment)};
\node[font=\footnotesize,text=omThm] at (5.2,-1.0) [align=center]{collatéral d'environ 102~\%~; les frais remontent vers la gauche,\\répartis à chaque étape};
\end{tikzpicture}

\omcaption{La chaîne de prêt. L'acheteur situé au bout est un propriétaire à part entière,
habilité à voter, à percevoir les dividendes et à prêter de nouveau les actions. Le détenteur
d'origine détient une créance sur son emprunteur, reçoit un paiement en remplacement de chaque
dividende, et peut exiger leur \omterm{def:m1:stock-loan-and-short-selling:recall}{rappel}.}\label{fig:m1:stock-loan-and-short-selling:chain}
\end{omfigure}

\begin{definition}[Encours vendeur à découvert et jours de couverture]\label{def:m1:stock-loan-and-short-selling:si}
L'\emph{encours vendeur à découvert}\index{encours vendeur à découvert} d'une action est le nombre de
ses actions vendues à découvert et non encore couvertes, souvent exprimé en pourcentage
du flottant. Les \emph{jours de couverture}\index{jours de couverture} se calculent en divisant l'encours vendeur
par le volume quotidien moyen.
\end{definition}

Comme l'acheteur d'une action vendue à découvert peut la prêter à nouveau, la même action
peut être vendue à découvert plus d'une fois, et l'encours vendeur peut dépasser 100~\% du flottant
sans aucun défaut de livraison~: les détenteurs longs se contentent de détenir
à eux tous le flottant plus l'encours vendeur.

\section{Frais, rebates et specials}

\begin{definition}[Collatéral général et special]\label{def:m1:stock-loan-and-short-selling:special}
Un titre est \emph{collatéral général}\index{collatéral général} (GC) lorsque l'offre
dépasse tellement la demande qu'il se prête à des frais minimes, de quelques dixièmes de pour cent
par an. Il est \emph{special}\index{special} (ou difficile à emprunter) lorsque les frais sont
nettement plus élevés~; le \omterm{def:m1:financing:seclending}{taux de rebate} d'un special est souvent négatif~: l'emprunteur
ne perçoit aucun intérêt sur son collatéral en espèces et paie en plus.
\end{definition}

Les frais représentent l'écart entre le taux d'intérêt de marché et le rebate
(\cref{def:m1:financing:seclending}). Ils sont fixés prêt par prêt, peuvent être modifiés
par le prêteur n'importe quel jour, et s'accumulent sur la valeur de marché \emph{courante} des
actions~: quand une vente à découvert tourne mal, le loyer augmente avec le cours. Les frais sont
négociés de gré à gré~; aucun flux public ne les publie. Les desks s'appuient sur les indications de leurs \omterm{def:m1:financing:pb}{prime
brokers} et sur les fournisseurs de données qui mettent en commun les livres de prêts des
prêteurs et des courtiers.

\begin{omfigure}
\begin{tikzpicture}
\begin{axis}[width=9.5cm,height=4.8cm,
  xlabel={taux d'utilisation (\%)},ylabel={frais d'emprunt (\% par an)},
  tick label style={font=\small},label style={font=\small},
  xmin=0,xmax=100,ymin=0,ymax=85,grid=major,grid style={gray!25},table/col sep=comma]
\addplot[omDef,very thick] table[x=utilisation_pct,y=fee_pct]{figdata/markets-1/16-stock-loan-and-short-selling/fee.csv};
\node[font=\footnotesize,anchor=south] at (axis cs:38,4) {collatéral général~: 0.3~\%};
\node[font=\footnotesize,anchor=east] at (axis cs:91,55) {special};
\end{axis}
\end{tikzpicture}

\omcaption{La courbe de frais illustrative du chapitre~: plate tant que l'offre est abondante,
convexe dès que les trois quarts de cette offre sont prêtés. Les courbes réelles varient selon le titre et
selon le jour~; la forme, un long palier plat suivi d'une fin abrupte, constitue l'élément
robuste.}\label{fig:m1:stock-loan-and-short-selling:fee}
\end{omfigure}

\begin{proposition}[Période de détention au seuil de rentabilité]\label{prop:m1:stock-loan-and-short-selling:breakeven}
Un vendeur à découvert anticipe une baisse du cours d'une fraction $g$. Avec des \omterm{def:m1:financing:seclending}{frais d'emprunt}
$\phi$ par an, les frais absorbent la totalité du gain attendu au bout de
\[
T^* \approx \frac{g}{\phi}\ \text{ans},
\]
au premier ordre (frais calculés sur la valeur initiale, intérêts sur le produit de la
vente négligés).
\end{proposition}

\begin{example}[Juste, et déficitaire]\label{ex:m1:stock-loan-and-short-selling:right}
Une thèse anticipant une baisse de 25~\% donne $T^* = 225$ jours à des frais de 40~\%, et 112 jours
à 80~\%. La \Cref{fig:m1:stock-loan-and-short-selling:carry} suit une vente à découvert de
100\,000 actions à 40~\$ sur une année simulée où le \omterm{def:m1:stock-loan-and-short-selling:utilisation}{taux d'utilisation} passe de
60~\% à 99~\%. Le titre monte d'abord à 53.52~\$ au jour 94, soit une perte de 1.36
million de dollars. Au jour 190, le titre est à 31.20~\$ et la position est
en avance de 789\,000~\$. Au jour 250, le cours est remonté à 38.36~\$, toujours sous
le prix d'entrée, et l'emprunt a coûté 399\,000~\$, dont les trois quarts
au cours des soixante derniers jours~: la transaction a perdu 235\,000~\$.
\end{example}

\begin{omfigure}
\begin{tikzpicture}
\begin{axis}[width=11cm,height=5.2cm,
  xlabel={jour},ylabel={milliers de dollars},
  tick label style={font=\small},label style={font=\small},
  xmin=1,xmax=250,ymin=-1500,ymax=1100,grid=major,grid style={gray!25},
  legend columns=3,legend style={font=\footnotesize,at={(0.5,-0.3)},anchor=north,draw=none,fill=none,/tikz/every even column/.append style={column sep=8pt}},
  table/col sep=comma]
\addplot[omExo,thick] table[x=day,y=price_pnl_k]{figdata/markets-1/16-stock-loan-and-short-selling/carry.csv};
\addplot[omThm,very thick,dashed] table[x=day,y=cost_k]{figdata/markets-1/16-stock-loan-and-short-selling/carry.csv};
\addplot[omDef,very thick] table[x=day,y=net_k]{figdata/markets-1/16-stock-loan-and-short-selling/carry.csv};
\legend{P\&L du cours,coût d'emprunt cumulé,net}
\end{axis}
\end{tikzpicture}

\omcaption{Une vente à découvert de 100\,000 actions à 40~\$ alors que le titre devient \omterm{def:m1:stock-loan-and-short-selling:special}{special}.
Le coût d'emprunt est négligeable pendant un semestre, puis absorbe tout.
Données~: la simulation du tutoriel.}\label{fig:m1:stock-loan-and-short-selling:carry}
\end{omfigure}

\section{Rappels, buy-ins et règles}

\begin{definition}[Rappel et buy-in]\label{def:m1:stock-loan-and-short-selling:recall}
Un \emph{rappel}\index{rappel} est la demande, adressée par le prêteur, de restitution des titres
prêtés, généralement parce qu'il les a vendus ou qu'il souhaite voter. Si le
courtier de l'emprunteur ne peut pas remplacer le prêt auprès d'un autre prêteur et
l'emprunteur ne restitue pas les actions à temps, le prêteur (ou le courtier)
exécute un \emph{buy-in}\index{buy-in}~: il achète les actions sur le marché aux
frais de l'emprunteur.
\end{definition}

\begin{definition}[Vente à découvert à nu]\label{def:m1:stock-loan-and-short-selling:naked}
Une \emph{vente à découvert à nu}\index{vente à découvert à nu} est une vente à découvert réalisée
sans avoir emprunté, ou pris des dispositions pour emprunter, le titre à temps pour
le règlement-livraison, de sorte que le vendeur ne parvient pas à livrer
(\cref{ch:m1:clearing-and-settlement}).
\end{definition}

\begin{dated}{2026-09}{Les règles de part et d'autre de l'Atlantique}\label{dat:m1:stock-loan-and-short-selling:rules}
\textbf{États-Unis, Regulation SHO.} \emph{Locate} (Rule 203(b)(1))~: avant
de réaliser une vente à découvert, un \omterm{def:m1:the-sell-side:sell-side}{broker-dealer} doit avoir des motifs raisonnables de
croire que le titre peut être emprunté et livré à la date d'échéance.
\emph{Close-out} (Rule 204)~: un défaut de livraison issu d'une vente à découvert doit être clôturé avant
l'ouverture de la séance régulière le jour de règlement-livraison qui suit la
date de règlement-livraison~; les défauts issus des ventes longues et de la tenue de marché de bonne foi disposent
jusqu'au troisième jour de règlement-livraison qui suit. \emph{Threshold securities}~:
défauts portant sur au moins 10\,000 actions et sur 0.5~\% des actions en circulation pendant cinq
jours de règlement-livraison consécutifs~; les défauts persistant pendant 13 jours
consécutifs de règlement-livraison doivent être clôturés sans délai. \emph{Price test} (Rule 201)~: après
baisse de 10~\% en une séance, les ventes à découvert sont limitées aux prix supérieurs au
meilleur cours acheteur national, pour le reste de cette séance et la suivante.
Le résumé publié par le régulateur indique que la \omterm{def:m1:stock-loan-and-short-selling:naked}{vente à découvert à nu} n'est pas en soi une infraction.

\textbf{Union européenne, Short Selling Regulation.} Les positions vendeuses nettes sur
les actions sont notifiées à l'autorité nationale à partir de 0.1~\% du capital
social émis, puis à chaque palier supplémentaire de 0.1~\%, et publiées à partir de 0.5~\%.
Les ventes à découvert non couvertes d'actions et de dette souveraine sont interdites, et les
autorités nationales peuvent imposer des interdictions temporaires en cas de crise.
\end{dated}

Les deux régimes prévoient une exception pour la tenue de marché~: aux États-Unis, un
\omterm{def:m1:the-sell-side:sell-side}{broker-dealer} exerçant une tenue de marché de bonne foi n'est pas tenu d'effectuer un locate
avant de vendre, car une société tenue de coter les deux côtés en permanence ne peut pas
consulter un desk de \omterm{def:m1:financing:seclending}{prêt de titres} avant chaque vente~; l'exception ne couvre ni la
vente spéculative ni la cotation unilatérale. L'exception et le délai de clôture plus long font
partie de l'économie des \omterm{def:m1:what-a-trading-firm-does:market-maker}{teneurs de marché} décrite au
\cref{ch:m1:what-a-trading-firm-does}.

\section{Short squeezes}

\begin{definition}[Short squeeze]\label{def:m1:stock-loan-and-short-selling:squeeze}
Un \emph{short squeeze}\index{short squeeze} est une hausse du cours amplifiée par
les achats des vendeurs à découvert qui se couvrent à cause de pertes, d'\omterm{def:m1:financing:margincall}{appels de marge},
de \omterm{def:m1:stock-loan-and-short-selling:recall}{rappels} ou de limites de risque, chaque achat poussant le cours encore plus haut
au détriment de ceux qui restent.
\end{definition}

\begin{proposition}[La cascade des couvertures]\label{prop:m1:stock-loan-and-short-selling:cascade}
Soient les seuils de couverture des vendeurs à découvert, exprimés en rendements par rapport
au cours du jour, répartis uniformément entre $a$ et $b$, et soit $k$ l'impact sur le cours, en
rendement, de la couverture de tous ces vendeurs. Après un choc exogène $s > a$, le cours
se fixe au rendement
\[
x \;=\; \frac{s - ka/(b-a)}{1 - k/(b-a)} \qquad\text{si } k < b-a \text{ et } x < b,
\]
et à $x = s + k$, avec tous les vendeurs couverts, dans le cas contraire. Le multiplicateur
$1/(1-k/(b-a))$ n'est pas borné lorsque $k \to b-a$.
\end{proposition}
\begin{proof}
La fraction couverte au rendement $x$ est $F(x) = (x-a)/(b-a)$ sur $[a,b]$, et le
cours vérifie $x = s + kF(x)$. On résout l'équation linéaire. Si $k \ge b-a$
chaque couverture en déclenche au moins une autre et l'itération $x \leftarrow s +
kF(x)$ ne s'arrête qu'en $F = 1$.
\end{proof}

\begin{omfigure}
\begin{tikzpicture}
\begin{axis}[width=10.5cm,height=5.4cm,
  xlabel={choc initial (\%)},ylabel={hausse finale (\%)},
  tick label style={font=\small},label style={font=\small},
  xmin=0,xmax=40,ymin=0,ymax=105,grid=major,grid style={gray!25},
  legend columns=4,legend style={font=\footnotesize,at={(0.5,-0.3)},anchor=north,draw=none,fill=none,/tikz/every even column/.append style={column sep=6pt}},
  table/col sep=comma]
\addplot[omExo,thick,dotted] coordinates {(0,0) (40,40)};
\addplot[omDef,very thick] table[x=shock_pct,y=k20]{figdata/markets-1/16-stock-loan-and-short-selling/squeeze.csv};
\addplot[omThm,very thick,dashed] table[x=shock_pct,y=k40]{figdata/markets-1/16-stock-loan-and-short-selling/squeeze.csv};
\addplot[black,very thick,dashdotted] table[x=shock_pct,y=k60]{figdata/markets-1/16-stock-loan-and-short-selling/squeeze.csv};
\legend{sans vendeurs à découvert,$k = 20\%$,$k = 40\%$,$k = 60\%$}
\end{axis}
\end{tikzpicture}

\omcaption{La cascade de couverture avec des seuils compris entre $+10\%$ et $+60\%$.
Avec $k = 20\%$, un choc est amplifié par $5/3$~; avec $k = 40\%$, par 5, jusqu'à ce qu'un
choc de 20~\% ait conduit tout le monde à couvrir~; avec $k = 60\% > b - a$, la première
couverture entraîne toutes les autres. Données~:
\cref{prop:m1:stock-loan-and-short-selling:cascade}.}\label{fig:m1:stock-loan-and-short-selling:squeeze}
\end{omfigure}

L'impact $k$ croît avec l'\omterm{def:m1:stock-loan-and-short-selling:si}{encours vendeur à découvert} rapporté aux actions qui
peuvent réellement être achetées~: les \omterm{def:m1:stock-loan-and-short-selling:si}{jours de couverture} et l'encours vendeur rapporté au \emph{flottant}
libre, sont les deux chiffres à surveiller. Deux épisodes marquent les extrêmes.

\begin{example}[Wolfsburg, octobre 2008]\label{ex:m1:stock-loan-and-short-selling:vw}
Le dimanche 26 octobre 2008, un constructeur automobile annonça qu'il détenait 42.6~\%
des actions ordinaires d'un autre constructeur, ainsi que des options réglées en espèces portant sur 31.5~\% supplémentaires. Un
État allemand en détenait 20~\%. Environ 13~\% des actions avaient été vendues à découvert, et 5.9~\%
restaient dans le flottant. Le cours, de 211~\euro le vendredi, dépassa 1\,005~\euro
le mardi 28 octobre, faisant brièvement de la société la plus valorisée parmi les
sociétés cotées au monde. Les options étaient l'enjeu~: étant réglées en espèces, elles
échappaient aux règles de transparence en vigueur à l'époque, et les actions que leurs
émetteurs détenaient en couverture n'étaient pas à vendre.
\end{example}

\begin{example}[janvier 2021]\label{ex:m1:stock-loan-and-short-selling:gme}
L'\omterm{def:m1:stock-loan-and-short-selling:si}{encours vendeur à découvert} d'un distributeur américain de jeux vidéo a atteint 122.97~\% de son flottant en
janvier 2021. Du plus bas intrajournalier du 8 janvier au plus haut intrajournalier de 483~\$
le 28 janvier, le cours a progressé d'environ 2\,700~\%, puis a chuté de plus de 86~\%
d'ici à la fin de la première semaine de février. Le rapport du personnel du régulateur a établi que
les rachats des vendeurs à découvert coïncidaient avec certaines phases de la hausse, mais n'étaient qu'une faible
fraction des achats totaux, et a conclu que le sentiment positif, et non les rachats
de couverture, avait soutenu la hausse~: un \omterm{def:m1:stock-loan-and-short-selling:squeeze}{short squeeze} faisait partie de l'histoire, mais non
de l'intégralité.
\end{example}

\section{Dividendes et fiscalité}

L'emprunteur d'une action doit au prêteur un \emph{dividende reconstitué} pour
chaque dividende versé pendant la durée du prêt. Le prêteur est économiquement indemne~; mais il
n'est plus le propriétaire légal à la \omterm{def:m1:shares-corporate-actions-indices:dividend}{date d'enregistrement}~; lorsqu'une retenue à la source
s'applique, la situation fiscale du propriétaire, de l'emprunteur et de l'acheteur peut différer. Prêter
des actions à cheval sur la \omterm{def:m1:shares-corporate-actions-indices:dividend}{date de détachement} du dividende à un détenteur dont la situation fiscale est plus favorable, et
en partager le gain, est une opération ancienne et, dans de nombreuses juridictions, contestée.

Sa dérive criminelle fut le \emph{cum-ex}~: des ventes à découvert organisées autour de la
\omterm{def:m1:shares-corporate-actions-indices:dividend}{date de détachement}, de sorte que la retenue à la source, acquittée une seule fois, semblait avoir été
versée à deux propriétaires, qui en demandèrent tous deux le remboursement. Le 28 juillet 2021, la
Cour fédérale de justice allemande a confirmé que réclamer le remboursement d'impôts qui n'ont jamais
été versés constitue une fraude fiscale, une infraction pénale. Pour un desk, la leçon est
étroite et absolue~: une opération sur dividende dont le profit est un remboursement
d'impôt relève d'une question juridique avant d'être une question de trading.

\section{Tutoriel~: le coût d'une vente à découvert lorsque l'action devient special}\label{tut:m1:stock-loan-and-short-selling:cost}

\begin{tutorial}
\textbf{Objectif.} Porter une vente à découvert pendant une année où l'emprunt se tend, et
reproduire la cascade de couverture. \textbf{État final~:} les trois figures de données
de ce chapitre.
\begin{tutsteps}
\item \textbf{Une courbe de frais.} Plate, puis convexe.
\omcode{code/markets-1/16-stock-loan-and-short-selling/python/short_cost.py}{5}{10}{Frais d'emprunt illustratif en fonction du taux d'utilisation.}\label{lst:m1:stock-loan-and-short-selling:fee}
\item \textbf{Carry.} Les frais s'accumulent sur la valeur de marché du jour, en actual/360.
\omcode{code/markets-1/16-stock-loan-and-short-selling/python/short_cost.py}{13}{17}{Résultat P\&L de cours et coût cumulé d'emprunt d'une vente à découvert.}\label{lst:m1:stock-loan-and-short-selling:carry}
\item \textbf{La cascade.} Une forme fermée et une itération qui doivent coïncider.
\omcode{code/markets-1/16-stock-loan-and-short-selling/python/short_cost.py}{29}{37}{L'équilibre de la cascade de couverture.}\label{lst:m1:stock-loan-and-short-selling:squeeze}
\end{tutsteps}
\textbf{Pour aller plus loin.} Faire réagir le \omterm{def:m1:stock-loan-and-short-selling:utilisation}{taux d'utilisation} au cours (les ventes
à découvert s'accroissent quand l'action monte sans nouvelle, les prêteurs vendent sur la hausse et rappellent).
Remplacer les seuils uniformes par une distribution concentrée près de
pertes à niveaux ronds, et observer ce qu'il advient du multiplicateur.
\end{tutorial}

\section{Construction~: le carnet d'emprunts}\label{bld:m1:stock-loan-and-short-selling:borrow}

\begin{build}
\bfield{Rôle} Le calculateur de financement du \cref{ch:m1:financing} prend
en entrée les \omterm{def:m1:financing:seclending}{frais d'emprunt}. Ce composant les produit, répond aux demandes de locate
et indique à la société le \omterm{def:m1:stock-loan-and-short-selling:recall}{buy-in} qu'elle doit exécuter lorsqu'un prêteur rappelle ses titres.
\bfield{Interface} \texttt{FeeCurve(gc, kink, max\_fee).fee(utilisation)};
\texttt{BorrowBook.locate(sym, qty)}, \texttt{borrow}, \texttt{give\_back},
\texttt{recall(sym, qty)} renvoyant la quantité de \omterm{def:m1:stock-loan-and-short-selling:recall}{buy-in} de la société,
\texttt{utilisation(sym)} et \texttt{borrow\_fees()}, un dictionnaire accepté
tel quel par \texttt{firm\_financing.accrue}.
\bfield{Règles} L'emprunt propre de la société compte dans le \omterm{def:m1:stock-loan-and-short-selling:utilisation}{taux d'utilisation}~: une grande vente
à découvert modifie ses propres frais. Un \omterm{def:m1:stock-loan-and-short-selling:recall}{rappel} qui laisse les prêts au-dessus de l'offre est réparti au prorata
entre la société et le reste du marché, en arrondissant au détriment de la société. Une
demande de locate pour un symbole inconnu renvoie zéro, jamais une erreur~: pas de locate, pas de
vente à découvert.
\bfield{Tests d'acceptation} \texttt{code/firm/borrow/tests/}~: demandes de locate et niveaux de
frais~; l'emprunt propre de la société modifiant les frais~; un \omterm{def:m1:stock-loan-and-short-selling:recall}{rappel} avec \omterm{def:m1:stock-loan-and-short-selling:recall}{buy-in} au prorata~;
le passage à la comptabilisation des frais de financement.
\bfield{Pour aller plus loin} Prêts à terme à frais fixes contre prêts overnight à
frais variables~: le choix auquel tout desk de \omterm{def:m1:financing:seclending}{prêt de titres} est confronté avant un événement.
\end{build}

\begin{omsources}
\item US Securities and Exchange Commission, \emph{Key Points About
Regulation SHO}, publication à destination des investisseurs.
\item US Securities and Exchange Commission, \emph{Staff Report on Equity and
Options Market Structure Conditions in Early 2021}, 14 octobre 2021.
\item European Securities and Markets Authority, \emph{Short selling}, page
d'orientation~; Regulation (EU) No 236/2012.
\item F.~Allen, M.~Haas, E.~Nowak et A.~Tengulov, «~Market efficiency and
limits to arbitrage: evidence from the Volkswagen short squeeze~»,
\emph{Journal of Financial Economics} 142 (2021)~; «~The case of Volkswagen~»,
\emph{The Hedge Fund Journal}.
\item Bundesgerichtshof, communiqué de presse 146/2021 sur l'arrêt 1 StR 519/20 du
28 juillet 2021.
\item D.~Duffie, N.~Gârleanu et L.~H.~Pedersen, «~Securities lending,
shorting, and pricing~», \emph{Journal of Financial Economics} 66 (2002).
\end{omsources}

\section{Exercices}

\begin{exercise}[$\star$]\label{exo:m1:stock-loan-and-short-selling:1}
Le taux au jour le jour est de 4.3~\% et un \omterm{def:m1:financing:pb}{prime broker} cote un \omterm{def:m1:financing:seclending}{taux de rebate} de $-11.7\%$ sur
un titre. Donner les \omterm{def:m1:financing:seclending}{frais d'emprunt} et le coût quotidien d'une vente à découvert de 3~millions de dollars.
\end{exercise}

\begin{exercise}[$\star$]\label{exo:m1:stock-loan-and-short-selling:2}
Un titre dispose de 12 millions d'actions d'\omterm{def:m1:stock-loan-and-short-selling:utilisation}{offre prêtable}, dont 10.2 millions
sont en prêt. Donner le \omterm{def:m1:stock-loan-and-short-selling:utilisation}{taux d'utilisation} et les frais lus sur la courbe du chapitre.
\end{exercise}

\begin{exercise}[$\star$]\label{exo:m1:stock-loan-and-short-selling:3}
Un titre a un flottant de 60 millions d'actions, un \omterm{def:m1:stock-loan-and-short-selling:si}{encours vendeur à découvert} de 18 millions et
un volume moyen de 2.4 millions par jour. Donner l'\omterm{def:m1:stock-loan-and-short-selling:si}{encours vendeur à découvert} en
pourcentage du flottant et les \omterm{def:m1:stock-loan-and-short-selling:si}{jours de couverture}. Lequel des deux compte le plus pour
un vendeur à découvert inquiet d'un \omterm{def:m1:stock-loan-and-short-selling:squeeze}{short squeeze}, et pourquoi~?
\end{exercise}

\begin{exercise}[$\star\star$]\label{exo:m1:stock-loan-and-short-selling:4}
Un analyste anticipe une baisse de 25~\%. Donner la durée de détention au point mort à des frais
de 40~\% et de 80~\%. Elle s'attend à ce que la baisse survienne «~à un moment dans les deux
prochaines années~» avec une probabilité uniforme. Estimer à peu près la part de ses scénarios
qui rapportent de l'argent à 40~\%~?
\end{exercise}

\begin{exercise}[$\star\star$]\label{exo:m1:stock-loan-and-short-selling:5}
Une société a un flottant de 100 actions, toutes détenues par A, qui les prête. Construire
une suite de prêts et de ventes après laquelle l'\omterm{def:m1:stock-loan-and-short-selling:si}{encours vendeur à découvert} atteint 123~\% du
flottant sans qu'aucune livraison n'ait échoué. Combien d'actions détiennent les longs
au total~? Qui vote~?
\end{exercise}

\begin{exercise}[$\star\star$]\label{exo:m1:stock-loan-and-short-selling:6}
Une vente à découvert exécutée le lundi 14 septembre 2026 ne se règle pas. En utilisant
l'\cref{dat:m1:stock-loan-and-short-selling:rules} et un règlement-livraison à un jour,
donner la date limite de clôture forcée. Quelle serait cette date pour une vente de position
longue ayant échoué~?
\end{exercise}

\begin{exercise}[$\star\star\star$]\label{exo:m1:stock-loan-and-short-selling:7}
\emph{Programmation.} À partir de \texttt{carry.csv} (ou en exécutant le script de la figure)
donner le P\&L net à son maximum avec son jour, le P\&L de prix final, le
coût cumulé final d'emprunt, et la part de ce coût engagée après le jour
190.
\end{exercise}

\begin{exercise}[$\star\star\star$]\label{exo:m1:stock-loan-and-short-selling:8}
\emph{Repérer le défaut.} «~L'\omterm{def:m1:stock-loan-and-short-selling:si}{encours vendeur à découvert} représente 140~\% du flottant. On a vendu plus
d'actions qu'il n'en existe. Cela prouve une \omterm{def:m1:stock-loan-and-short-selling:naked}{vente à découvert à nu} illégale, avec des
actions contrefaites.~» Expliquer pourquoi l'inférence échoue, et dire quelles données \emph{constitueraient}
une preuve de \omterm{def:m1:stock-loan-and-short-selling:naked}{vente à découvert à nu}.
\end{exercise}

\section{Problème~: Le short squeeze}

\begin{problem}[{Problème du week-end --- un short surpeuplé face à une bonne nouvelle}]\label{pb:m1:stock-loan-and-short-selling:1}
Un fonds est vendeur à découvert de 2 millions d'actions de Z à 20~\$. Z a un flottant de 50 millions
d'actions, un \omterm{def:m1:stock-loan-and-short-selling:si}{encours vendeur à découvert} de 20 millions et un volume moyen de 4 millions par
jour. Le fonds anticipe une baisse de 30~\%. Les \omterm{def:m1:financing:seclending}{frais d'emprunt} sont de 12~\%.

\medskip\noindent\textbf{Partie I --- Le carry.}
\begin{enumerate}
\item Donner l'\omterm{def:m1:stock-loan-and-short-selling:si}{encours vendeur à découvert} en pourcentage du flottant, les \omterm{def:m1:stock-loan-and-short-selling:si}{jours de couverture},
et la part du fonds dans l'\omterm{def:m1:stock-loan-and-short-selling:si}{encours vendeur à découvert}.
\item Donner le coût quotidien de l'emprunt.
\item Donner la durée de détention au point mort.
\item Les frais sont révisés à 60~\%. Donner le nouveau coût quotidien et le nouveau point mort.
\item La durée moyenne de détention du fonds pour de telles opérations est de neuf mois. Que
implique la question 4~?
\end{enumerate}

\medskip\noindent\textbf{Partie II --- La nouvelle.}
Z grimpe de 25~\% sur une rumeur d'OPA. Les seuils de rachat des vendeurs à découvert sont
uniformes entre $+10\%$ et $+60\%$.
\begin{enumerate}[resume]
\item Avec $k = 20\%$, à quel niveau se stabilise le cours, et quelle part des
vendeurs à découvert a racheté~?
\item Avec $k = 40\%$~?
\item Pour $k = 40\%$, déterminer le plus petit choc après lequel tous les vendeurs à découvert rachètent.
\item Donner la perte de \omterm{def:m1:pnl-and-positions:mtm}{valorisation au prix de marché} du fonds au cours obtenu à la question 7.
\item Son \omterm{def:m1:financing:pb}{prime broker} relève la marge sur la vente à découvert de 30~\% à 100~\% de la
valeur de marché. Donner la marge supplémentaire à ce cours.
\end{enumerate}

\medskip\noindent\textbf{Partie III --- Le \omterm{def:m1:stock-loan-and-short-selling:recall}{rappel}.}
\begin{enumerate}[resume]
\item Un prêteur rappelle les titres~; le courtier ne peut pas remplacer 500\,000 des
actions du fonds et les rachète par \omterm{def:m1:stock-loan-and-short-selling:recall}{buy-in} à 33~\$. Donner la perte réalisée.
\item Pourquoi les \omterm{def:m1:stock-loan-and-short-selling:recall}{rappels} se concentrent-ils précisément quand le cours a augmenté~?
\item Un mois plus tôt, le fonds avait le choix d'un emprunt à terme de trois mois
à taux fixe de 20~\%. Donner son coût sur trois mois et l'argument en faveur de
ce paiement.
\item Proposer une position sur options cotées avec la même anticipation que la vente à découvert
sans emprunt. Quel en est le coût à la place, et où sont passés les \omterm{def:m1:financing:seclending}{frais d'emprunt}~?
\end{enumerate}

\medskip\noindent\textbf{Partie IV --- Jugement.}
\begin{enumerate}[resume]
\item Lesquels des chiffres de la question 1 auraient alerté le fonds~?
\item Un collègue affirme que le \omterm{def:m1:stock-loan-and-short-selling:squeeze}{short squeeze} «~prouve une manipulation~». En s'appuyant sur le
rapport du personnel figurant dans l'\cref{ex:m1:stock-loan-and-short-selling:gme}, expliquer pourquoi cela
exige une preuve propre.
\item La rumeur s'avère fausse et Z chute à 13~\$ en un trimestre. Le
fonds avait-il raison~?
\item Comment dimensionner les positions vendeuses, compte tenu des parties II et III~?
\item Énoncer le \emph{résultat nommé}~: la durée de détention au point mort de la
vente à découvert après la révision des frais.
\item En une phrase~: que loue un vendeur à découvert, et à qui~?
\end{enumerate}
\end{problem}

\section{Questions d'entretien}

\begin{interviewq}[$\star$ \iqroles{trader, chercheur}]\label{iq:m1:stock-loan-and-short-selling:1}
Décrivez-moi ce qui se passe, sur le plan opérationnel, lorsqu'un fonds vend une action à découvert~?
\end{interviewq}

\begin{interviewq}[$\star$ \iqroles{trader, chercheur}]\label{iq:m1:stock-loan-and-short-selling:2}
Comment l'\omterm{def:m1:stock-loan-and-short-selling:si}{encours vendeur à découvert} peut-il dépasser 100~\% du flottant~?
\end{interviewq}

\begin{interviewq}[$\star\star$ \iqroles{chercheur, ingénieur ML}]\label{iq:m1:stock-loan-and-short-selling:3}
Votre backtest vend à découvert le décile inférieur d'un signal et gagne 9~\% par an sur
la jambe vendeuse. Qu'est-ce qui manque~?
\end{interviewq}

\begin{interviewq}[$\star\star$ \iqroles{trader, chercheur}]\label{iq:m1:stock-loan-and-short-selling:4}
Les \omterm{def:m1:financing:seclending}{frais d'emprunt} d'une action sont passés de 2~\% à 45~\% en une semaine. Qu'est-ce que cela
vous apprend, et ce qu'il ne dit pas~?
\end{interviewq}

\begin{interviewq}[$\star\star$ \iqroles{trader, banque}]\label{iq:m1:stock-loan-and-short-selling:5}
Pourquoi les \omterm{def:m1:what-a-trading-firm-does:market-maker}{teneurs de marché} sont-ils dispensés de l'obligation de locate, et qu'est-ce qui
les empêche d'en abuser~?
\end{interviewq}

\begin{interviewq}[$\star\star\star$ \iqroles{chercheur, trader}]\label{iq:m1:stock-loan-and-short-selling:6}
Les options de vente sur une action difficile à emprunter paraissent chères par rapport aux calls~: la
parité call-put semble violée. Y a-t-il un arbitrage~?
\end{interviewq}
